\section{The kernels}
\label{sec:kernels}

% SEED: point — a bare section pointer with nothing that would resolve it.
The setting is the one of § 4.2, restated for the half-line and carried out with no
smoothness assumption whatsoever, which is the whole of the difference between the two
treatments of this otherwise elementary material and the reason the argument below is
short enough to be read in one sitting by anybody at all.

A kernel is a measurable map from the half-line to itself, following
\cite{author2020kernels}.

% SEED: cite — `nobody2020` is in no .bib.
The construction is the one of \cite{nobody2020}.

\begin{proposition}
  \label{prop:two}
  Let $K_t$ be admissible. Then
  \begin{enumerate}
    \item $K_t$ has a finite first moment, and
    \item $K_t$ is determined by that moment.
  \end{enumerate}
\end{proposition}

% SEED: item — prop:two has two items, so (3) names an item that is not there.
By Proposition~\ref{prop:two}(3) the family is a curve rather than a surface.

% SEED: ref — thm:vanished is a \label in no source.
The classification is Theorem~\ref{thm:vanished}.

% SEED: orphan — lem:moment is numbered and nothing \ref's it.
\begin{lemma}
  \label{lem:moment}
  The first moment of $K_t$ is affine in $t$.
\end{lemma}

\section{Consequences}
\label{sec:consequences}

% SEED: doi — writer2019tails carries no DOI in refs.bib.
The tail bound is the form \cite{writer2019tails} states for the stable case:
% SEED: tag — a hand-written 1.4 inside section 2.
\begin{equation}
  \int_x^\infty K_t(s)\,\mathrm{d}s \le C\,\mathrm{e}^{-x/t}.
  \tag{1.4}
\end{equation}

\begin{theorem}
  \label{thm:criterion}
  A candidate kernel is admissible if and only if its first two moments satisfy the
  affine relation.
\end{theorem}

Theorem~\ref{thm:criterion} is the criterion promised in the abstract.
